\documentclass[11pt]{article}
\usepackage{amsmath,amssymb,amsthm,mathtools}
\usepackage[margin=1.1in]{geometry}
\usepackage{hyperref}

\theoremstyle{plain}
\newtheorem{theorem}{Theorem}[section]
\newtheorem{proposition}[theorem]{Proposition}
\newtheorem{lemma}[theorem]{Lemma}
\newtheorem{corollary}[theorem]{Corollary}
\newtheorem{conjecture}[theorem]{Conjecture}
\theoremstyle{definition}
\newtheorem{definition}[theorem]{Definition}
\newtheorem{remark}[theorem]{Remark}
\newtheorem{example}[theorem]{Example}

\newcommand{\Z}{\mathbb{Z}}
\newcommand{\PFtwo}{\mathrm{PF}_2}
\newcommand{\Bx}{\mathrm{Box}}
\DeclareMathOperator{\supp}{supp}

\title{Condition (A): log-concavity of cylindric Kostka numbers at $\ell=3$\\
{\large A reformulation valid for all $m$, a complete reduction at $m=2$,\\
and a proof in the single-region case}}
\author{Clio}
\date{30 September 2026}

\begin{document}
\maketitle

\begin{abstract}
Let $s^{c}_{\lambda/\mu}(x_1,x_2,x_3)$ be a cylindric skew Schur polynomial at level $n$
with $m$ beads, and let $k(a,b)=K^{c}_{(a,b,d-a-b)}$ be its cylindric Kostka numbers.
Condition~(A) is the assertion $k(a,b)^2\ge k(a+1,b)k(a-1,b)$; by the $S_3$-symmetry of
$s^c$ this is log-concavity along every root direction, and it is one of the two
conjuncts that would settle the Lorentzian property~(L3) at $\ell=3$.

We prove three things and record one gap.
(i) For \emph{every} $m$, condition~(A) is equivalent to the statement that a sum of
$\PFtwo$ sequences over a single slice of a box is log-concave
(Proposition~\ref{prop:reform}); the summands are exactly the $\ell=2$ cylindric Kostka
sequences, which are $\PFtwo$ by an already-proved theorem. This converts~(A) from a
statement about a three-step chain into a one-parameter statement.
(ii) At $m=2$ the slice is one-dimensional, and~(A) reduces \emph{exactly} to a
self-contained combinatorial lemma about a concave and a convex $1$-Lipschitz integer
profile (Proposition~\ref{prop:m2red}, Conjecture~\ref{conj:T}). The $1$-Lipschitz
hypothesis is sharp: dropping it produces counterexamples.
(iii) At $m=2$ the parameter range splits at two explicit breakpoints into at most four
regions, each of which contributes a $\PFtwo$ sequence, by two different mechanisms which
we prove (Theorem~\ref{thm:regions}). Consequently~(A) holds whenever a single region is
active --- an explicitly checkable criterion satisfied by $79.9\%$ of the slices in our
sweep (Corollary~\ref{cor:single}).

The remaining gap is precisely stated: the pairwise mixed inequality between two
\emph{distinct} regions (verified on $525$ of $525$ region pairs, unproved).

We emphasise what is \emph{not} claimed. Condition~(A) alone does not settle~(L3), even at
$\ell=3$: by \texttt{l3-det-reduction}, (L3) at $\ell=3$ is equivalent to
$e_2(M)\le 0$ \emph{and} $\det M\ge 0$, and (A) delivers only the first conjunct.
\end{abstract}

\section{Setting and statement}

\subsection{Cylindric shapes}
Fix a level $n\ge 2$ and a number of beads $m\ge 1$. A \emph{cylindric shape} is a tuple
$x=(x_1,\dots,x_m)\in\Z^m$ with $x_1<x_2<\dots<x_m<x_1+n$, extended cyclically by
$x_{i+m}=x_i+n$. For cylindric shapes $\nu,\rho$ we write $\nu\prec\rho$
(\emph{$\rho/\nu$ is a horizontal strip}) if
\[
  \nu_i\ \le\ \rho_i\ \le\ \nu_{i+1}-1\qquad (1\le i\le m),
\]
cyclically, and set $|\rho/\nu|=\sum_i(\rho_i-\nu_i)$. For a cylindric skew shape
$\lambda/\mu$ (so $\mu_i\le\lambda_i$ for all $i$) with $d=|\lambda/\mu|$, the cylindric
Kostka number $K^c_{(a_1,\dots,a_\ell)}$ counts chains
$\mu=\rho^0\prec\rho^1\prec\dots\prec\rho^{\ell}=\lambda$ with $|\rho^t/\rho^{t-1}|=a_t$,
and $s^c_{\lambda/\mu}(x_1,\dots,x_\ell)=\sum_\alpha K^c_\alpha x^\alpha$.

Throughout, $\ell=3$ and $k(a,b):=K^c_{(a,b,d-a-b)}$. Since $s^c$ is symmetric,
$k(a,b)=k(b,a)$.

\begin{conjecture}[Condition (A)]\label{conj:A}
For every cylindric skew shape $\lambda/\mu$ and all $a,b$,
$k(a,b)^2\ \ge\ k(a+1,b)\,k(a-1,b)$.
\end{conjecture}

\subsection{Inputs used}
We use the following, all recorded as \texttt{proved} in the registry
\texttt{proofs/registry/cylindric-lorentzian.json}.

\begin{itemize}
\item[(P1)] (\texttt{interlacing-model}) At $\ell=3$, $K^c_{(a,b,c)}$ counts pairs
$(\kappa,\sigma)\in\Z^m\times\Z^m$ with $\mu_i\le\kappa_i\le\mu_{i+1}-1$,
$\lambda_{i-1}+1\le\sigma_i\le\lambda_i$ and $\kappa_i\le\sigma_i\le\kappa_{i+1}-1$
(cyclically), graded by $a=S_\kappa-|\mu|$, $b=S_\sigma-S_\kappa$, $c=|\lambda|-S_\sigma$.
\item[(P2)] (\texttt{ell2-decoupling}) At $\ell=2$ the constraints on the single
intermediate shape decouple into one independent interval per bead.
\item[(P3)] (\texttt{pf2-convolution}) $\PFtwo$ --- log-concave with interval support --- is
closed under convolution.
\item[(P4)] (\texttt{ell2-all-m}) $\sum_a K^c_{\lambda/\mu,(a,d-a)}z^a$ has $\PFtwo$
coefficients, for every cylindric skew shape, every $n$, every $m$.
\end{itemize}

We write $\PFtwo$ for: nonnegative, support an interval of $\Z$, and
$g(s)^2\ge g(s-1)g(s+1)$ for all $s$.

\begin{lemma}[concave $\Rightarrow$ $\PFtwo$]\label{lem:conc}
Let $g:\Z\to\Z_{\ge0}$ have support an interval $J=[j_0,j_1]$ and satisfy
$2g(s)\ge g(s-1)+g(s+1)$ for all $j_0<s<j_1$. Then $g\in\PFtwo$.
\end{lemma}
\begin{proof}
For $j_0<s<j_1$, AM--GM gives
$g(s)\ge\tfrac12(g(s-1)+g(s+1))\ge\sqrt{g(s-1)g(s+1)}$, so $g(s)^2\ge g(s-1)g(s+1)$.
For $s=j_0$ we have $g(j_0-1)=0$, so the right side vanishes; likewise at $s=j_1$.
For $s\notin J$ at least one of $g(s-1),g(s+1)$ lies outside $J$ or $g(s)$ is on the
boundary; in every case $g(s-1)g(s+1)=0$ unless $s-1,s+1\in J$, which forces $s\in J$.
The support is an interval by hypothesis.
\end{proof}

\begin{lemma}[truncated concave]\label{lem:trunc}
Let $c:\Z\to\Z$ be concave (that is, $2c(s)\ge c(s-1)+c(s+1)$ for all $s$). Then
$g=\max(c,0)$ has interval support and is $\PFtwo$.
\end{lemma}
\begin{proof}
Put $J=\{c>0\}$.

\emph{$J$ is an interval.} A concave integer sequence is quasiconcave:
$c(s)\ge\min\bigl(c(r),c(t)\bigr)$ whenever $r\le s\le t$. (Induction on the span $t-r$:
for $r<s<t$ apply the statement to $[r,t-1]$ and to $[r+1,t]$ and compare the two slopes,
using that increments $c(\cdot+1)-c(\cdot)$ are non-increasing.) So $c(r)>0$ and $c(t)>0$
force $c(s)>0$ for every $s$ between them, which is exactly interval support.

\emph{Log-concavity.} Fix $s$. If $c(s-1)\le0$ then $g(s-1)=0$ and
$g(s-1)g(s+1)=0\le g(s)^2$; likewise if $c(s+1)\le0$. If $c(s-1)>0$ and $c(s+1)>0$ then
quasiconcavity gives $c(s)>0$ as well, so $g$ agrees with $c$ at all three points, and
$2c(s)\ge c(s-1)+c(s+1)$ with $c(s\pm1)\ge0$ gives
$c(s)^2\ge\tfrac14\bigl(c(s-1)+c(s+1)\bigr)^2\ge c(s-1)c(s+1)$ by AM--GM.

Note that this argument does \emph{not} pass through Lemma~\ref{lem:conc}. That lemma
requires the support to be a \emph{bounded} interval $J=[j_0,j_1]$, and $\max(c,0)$ need
not have bounded support: $c\equiv1$ is concave with $\{c>0\}=\Z$, so there is no $j_1$.
The proof above needs no boundedness.
\end{proof}

\begin{remark}\label{rem:truncLean}
Lemma~\ref{lem:trunc} is machine-checked, for $c:\Z\to\Z$, as
\verb|TworowD4Kernel.DiscreteConcavity.PFtwo_posPart| in
\verb|lean/tworow_d4_kernel/TworowD4Kernel/DiscreteConcavity.lean| (axioms: \verb|propext|,
\verb|Quot.sound| only); the quasiconcavity step is \verb|IntConcave.min_le| and the
counterexample to citing Lemma~\ref{lem:conc} is \verb|posPart_support_unbounded|. The
earlier version of this proof read ``at the two ends of the support the argument of
Lemma~\ref{lem:conc} applies verbatim'', and also allowed $c$ to take the value $-\infty$
off an interval. The first clause was a citation whose object is undefined in the
unbounded case; the second is not needed, since at the only place Lemma~\ref{lem:trunc} is
used --- $w_{\mathrm I}$ and $w_{\mathrm{III}}$ in the proof of
Proposition~\ref{prop:regI} --- the function $c$ is a genuine $\Z$-valued concave
sequence, being minus a maximum of finitely many affine functions.
\end{remark}

\section{A reformulation of (A), valid for every $m$}

\begin{proposition}\label{prop:reform}
Let $\Bx(\mu)=\prod_{i=1}^{m}[\mu_i,\mu_{i+1}-1]\subset\Z^m$ (with $\mu_{m+1}=\mu_1+n$),
and for $\nu\in\Bx(\mu)$ put
\[
  L_i(\nu)=\max(\nu_i,\lambda_{i-1}+1),\qquad
  R_i(\nu)=\min(\nu_{i+1}-1,\lambda_i),\qquad
  f_\nu \;=\; \mathop{\text{\Large$*$}}_{i=1}^{m}\ \mathbf 1_{[\,L_i(\nu)-\nu_i,\ R_i(\nu)-\nu_i\,]} .
\]
Then for all $a,b$,
\[
  k(a,b)\;=\;k(b,a)\;=\!\!\sum_{\substack{\nu\in\Bx(\mu)\\ S_\nu=|\mu|+b}}\!\! f_\nu(a),
\]
and each $f_\nu$ is $\PFtwo$. Consequently condition~(A) is equivalent to:
\emph{for every $b$, the sum of the $\PFtwo$ sequences $f_\nu$ over the box-slice
$\{\nu\in\Bx(\mu):S_\nu=|\mu|+b\}$ is log-concave.}
\end{proposition}

\begin{proof}
By (P1) with the roles of the two interior layers named $\nu,\kappa$, $k(b,a)$ counts pairs
with $|\nu/\mu|=b$, $|\kappa/\nu|=a$. The condition $\mu\prec\nu$ reads
$\mu_i\le\nu_i\le\mu_{i+1}-1$, one independent interval per bead, so the admissible $\nu$
form exactly the box $\Bx(\mu)$; no separate shape check is needed, because
$\nu_i\le\mu_{i+1}-1<\mu_{i+1}\le\nu_{i+1}$. Given $\nu$, the conditions $\nu\prec\kappa$
and $\kappa\prec\lambda$ read $\nu_i\le\kappa_i\le\nu_{i+1}-1$ and
$\lambda_{i-1}+1\le\kappa_i\le\lambda_i$, i.e.\ $\kappa_i\in[L_i(\nu),R_i(\nu)]$, again one
independent interval per bead; and $\kappa$ is automatically a cylindric shape because
$\kappa_i\le\nu_{i+1}-1<\nu_{i+1}\le\kappa_{i+1}$. Grading by
$a=\sum_i(\kappa_i-\nu_i)$ gives the convolution $f_\nu$. Each $f_\nu$ is a convolution of
interval indicators, hence $\PFtwo$ by (P3); this is (P2)+(P4) for the shape
$\lambda/\nu$. Finally $k(a,b)=k(b,a)$ by the symmetry of $s^c$.
\end{proof}

\begin{remark}
The point of Proposition~\ref{prop:reform} is that it removes one of the two interior
layers. Condition~(A) is \emph{not} a statement about a three-step chain; it is the
statement that a particular finite family of $\PFtwo$ sequences, indexed by the lattice
points of a box lying on one hyperplane, sums to a log-concave sequence.
\end{remark}

\subsection{What this rules out}

The following were tested against Proposition~\ref{prop:reform} and \emph{fail}; each is
recorded with the scope of its reason, since each is a route that looks open until it is
tried.

\begin{enumerate}
\item\textbf{Pairwise comparison inside a slice is not enough.} For sequences
$g_1,\dots,g_N\ge 0$, expanding $(\sum g)^2-(\sum g)(s-1)(\sum g)(s+1)$ and symmetrising
shows that $\sum g_j$ is log-concave as soon as each $g_j$ is log-concave and every pair
satisfies
\begin{equation}\label{eq:star}
  2g_j(s)g_k(s)\ \ge\ g_j(s-1)g_k(s+1)+g_j(s+1)g_k(s-1)\qquad\text{for all }s. \tag{$*$}
\end{equation}
This pairwise route \emph{fails} inside a slice. Witness: $n=6$, $m=2$, $\mu=(0,3)$,
$\lambda=(4,5)$, slice $S_\nu=5$, $\nu=(0,5)$ and $\rho=(2,3)$; then
$f_\nu=\mathbf 1_{[0,4]}$ and $f_\rho=\delta_2$, and at $s=1$ the left side
of~\eqref{eq:star} is $0$ while the right side is $1$. Over $n\le 6$, $m\ge2$, $d\le 7$,
$4$ of $9871$ pairs fail, and $244$ of $351$ distinct pairs are not $\mathrm{TP}_2$-comparable in
either orientation. \emph{Scope:} the slice family is a genuine antichain; the failing pair
is repaired only by the intervening element $\nu=(1,4)$ of the same slice. Any proof must
see the whole slice at once.

\item\textbf{The statement is not a general fact about difference-constraint polytopes.}
The set in (P1) is cut out by constraints $x_p-x_q\le c$ alone. But the corresponding
general assertion --- for $P=\{x\in\Z^{2m}: x_p-x_q\le c_{pq},\ \text{box}\}$, the fibre
count $v\mapsto\#\{x\in P:S_\kappa=u,\ S_\sigma=v\}$ is log-concave --- is \emph{false}:
$312$ failures out of $6474$ fibres in a random sweep with $m=2$; e.g.\ box
$0\le x\le(1,1,2,4)$ with constraints
$x_3-x_1\le1,\ x_3-x_1\le0,\ x_1-x_0\le3,\ x_1-x_0\le2,\ x_2-x_3\le1$ gives
$(2,3,1,1)$ at $u=1$. \emph{Scope:} the cylindric interlacing pattern, not mere
difference-constraint structure, is load-bearing.

\item\textbf{$f_b$ is not real-rooted}, so Newton's inequalities are unavailable: already
$f_b=1+x+x^2$ occurs ($n=3$, $\mu=(0)$, $\lambda=(2)$, $b=0$). $600$ of $1258$ slices are
not real-rooted.

\item\textbf{$f_b$ is not a convolution of interval indicators.} $91$ of $1176$ slice
polynomials are not products of cyclotomics; the smallest is $2+3x+2x^2$
($n=5$, $\mu=(0,2)$, $\lambda=(2,4)$), which is irreducible over $\mathbb Q$.

\item\textbf{Concavity is false} --- see Section~\ref{sec:conc}. This matters because the
proof of Proposition~\ref{prop:regII} below proves concavity, so it \emph{cannot} extend
to the general case.
\end{enumerate}

\section{The case $m=2$}

From here on $m=2$, $\mu=(\mu_1,\mu_2)$, $\lambda=(\lambda_1,\lambda_2)$, and we abbreviate
\[
  A=\lambda_2-n+1,\qquad B=\lambda_1,\qquad C=\lambda_1+1,\qquad D=\lambda_2,
\]
so that $C=B+1$ and $A=D-n+1$. Fix $b$ and put $u=|\mu|+b$. The slice
$\{\nu\in\Bx(\mu):S_\nu=u\}$ is parametrised by $t=\nu_1$, $\nu_2=u-t$, with
$t\in T=[T_-,T_+]$, $T_-=\max(\mu_1,u-\mu_1-n+1)$, $T_+=\min(\mu_2-1,u-\mu_2)$.
Given $t$, Proposition~\ref{prop:reform} says $\kappa_1$ and $\kappa_2$ range independently
over
\[
  \kappa_1\in[a_t,b_t]:=[\max(t,A),\ \min(u-1-t,B)],\qquad
  \kappa_2\in[c_t,d_t]:=[\max(u-t,C),\ \min(t+n-1,D)] .
\]
Writing $s=\kappa_1+\kappa_2=u+a$, we must show that
\begin{equation}\label{eq:G}
  G(s)\;=\;\sum_{t\in T}\bigl(\mathbf 1_{[a_t,b_t]}*\mathbf 1_{[c_t,d_t]}\bigr)(s)
\end{equation}
is $\PFtwo$.

\subsection{Exact reduction to a two-profile lemma}

\begin{proposition}\label{prop:m2red}
Define $\Phi:[A,B]\to\Z$ and $\Psi:[C,D]\to\Z$ by
\[
  \Phi(i)=\min(i,\ u-1-i,\ T_+),\qquad \Psi(j)=\max(u-j,\ j-n+1,\ T_-).
\]
Then $\Phi$ is concave and $\Psi$ is convex, both are $1$-Lipschitz, and
\begin{equation}\label{eq:PhiPsi}
  G(s)\;=\;\sum_{i+j=s}\bigl(\Phi(i)-\Psi(j)+1\bigr)_+ ,
\end{equation}
the sum being over $i\in[A,B]$, $j\in[C,D]$.
\end{proposition}
\begin{proof}
$\kappa_1=i$ is admissible for the parameter $t$ iff $A\le i\le B$, $t\le i$ and
$t\le u-1-i$; together with $t\in T$ this says $T_-\le t\le\Phi(i)$. Likewise $\kappa_2=j$
is admissible iff $C\le j\le D$ and $t\ge u-j$, $t\ge j-n+1$, i.e.\ $\Psi(j)\le t\le T_+$.
Hence the number of $t\in T$ for which the pair $(i,j)$ is counted is
$\#\{t:\Psi(j)\le t\le\Phi(i)\}=(\Phi(i)-\Psi(j)+1)_+$, giving~\eqref{eq:PhiPsi}. $\Phi$ is
a minimum of three affine functions of slope $1,-1,0$, hence concave and $1$-Lipschitz;
$\Psi$ is a maximum of three affine functions of slope $-1,1,0$, hence convex and
$1$-Lipschitz.
\end{proof}

This is an \emph{exact} reduction: nothing about cylindric shapes survives
in~\eqref{eq:PhiPsi} except the four numbers $A,B,C,D$ and the two numbers $u,n$, all
absorbed into $\Phi,\Psi$. So condition~(A) at $m=2$ is implied by:

\begin{conjecture}[Lemma T]\label{conj:T}
Let $\Phi$ be concave and $1$-Lipschitz on an integer interval, $\Psi$ convex and
$1$-Lipschitz on an integer interval (both extended by $\mp\infty$). Then
$s\mapsto\sum_{i+j=s}(\Phi(i)-\Psi(j)+1)_+$ is log-concave.
\end{conjecture}

\begin{remark}[the $1$-Lipschitz hypothesis is sharp]\label{rem:sharp}
Both weakenings fail.
\begin{itemize}
\item Without any slope bound: $\Phi=(-1,0,1)$, $\Psi=(-3,0,4)$ on $\{0,1,2\}$ give
$G=(3,4,6,2)$, and $4^2=16<3\cdot6=18$. ($40$ failures in $19133$ random trials.)
\item With slopes in $\{-2,\dots,2\}$: $\Phi=(3,5,5)$, $\Psi=(3,1,2,3,4)$ give
$G=(1,6,10,10,7,5,2)$, which fails at $s=4$: $7^2=49<10\cdot5=50$.
($117$ failures in $32900$ random trials.)
\item With slopes in $\{-1,0,1\}$: $0$ failures in $29251$ random trials, and $0$ failures
in an \emph{exhaustive} enumeration of all such profiles of length $\le6$ with base shifts
in $[-4,6]$ ($29000$ instances).
\end{itemize}
So the negative control bites, and the hypothesis is load-bearing rather than decorative.
\end{remark}

\subsection{The two breakpoints and the four regions}

The four endpoint functions in~\eqref{eq:G} are each a $\min$ or $\max$ of two affine
functions, and --- this is the structural point --- they share only \emph{two} breakpoints:
\[
  \tau_1:=A=D-n+1 \quad(\text{for }a_t\text{ and }d_t),
  \qquad
  \tau_2:=u-1-B=u-C \quad(\text{for }b_t\text{ and }c_t).
\]
Indeed $a_t=\max(t,A)$ switches at $t=A$ and $d_t=\min(t+n-1,D)$ switches at
$t=D-n+1=A$; while $b_t=\min(u-1-t,B)$ switches at $t=u-1-B$ and $c_t=\max(u-t,C)$
switches at $t=u-C=u-1-B$. Accordingly $T$ meets at most four cells:

\begin{center}
\begin{tabular}{llllll}
\hline
region & condition on $t$ & $a_t$ & $b_t$ & $c_t$ & $d_t$\\
\hline
I   & $t<\tau_1$, $t\le\tau_2$ & $A$ & $B$ & $u-t$ & $t+n-1$\\
II  & $\tau_1\le t\le\tau_2$   & $t$ & $B$ & $u-t$ & $D$\\
III & $t\ge\tau_1$, $t>\tau_2$ & $t$ & $u-1-t$ & $C$ & $D$\\
IV  & $t<\tau_1$, $t>\tau_2$   & $A$ & $u-1-t$ & $C$ & $t+n-1$\\
\hline
\end{tabular}
\end{center}

Write $G=G_{\mathrm I}+G_{\mathrm{II}}+G_{\mathrm{III}}+G_{\mathrm{IV}}$ for the
corresponding partial sums in~\eqref{eq:G}.

\begin{proposition}[regions I and III: a fixed factor]\label{prop:regI}
Let $T_{\mathrm I}=[p_0,p_1]$ be the (possibly empty) set of $t\in T$ in region I. Then
\[
  G_{\mathrm I}=\mathbf 1_{[A,B]}*w_{\mathrm I},\qquad
  w_{\mathrm I}(x)=\bigl(p_1-\max(u-x,\ x-n+1,\ p_0)+1\bigr)_+ ,
\]
and $G_{\mathrm I}\in\PFtwo$. Symmetrically, with $T_{\mathrm{III}}=[q_0,q_1]$,
\[
  G_{\mathrm{III}}=w_{\mathrm{III}}*\mathbf 1_{[C,D]},\qquad
  w_{\mathrm{III}}(x)=\bigl(\min(x,\ u-1-x,\ q_1)-q_0+1\bigr)_+ ,
\]
and $G_{\mathrm{III}}\in\PFtwo$.
\end{proposition}
\begin{proof}
In region I the first interval $[a_t,b_t]=[A,B]$ does not depend on $t$, so
\[
 G_{\mathrm I}=\sum_{t\in T_{\mathrm I}}\mathbf 1_{[A,B]}*\mathbf 1_{[u-t,\,t+n-1]}
  =\mathbf 1_{[A,B]}*\Bigl(\sum_{t\in T_{\mathrm I}}\mathbf 1_{[u-t,\,t+n-1]}\Bigr).
\]
Now $u-t\le x\le t+n-1$ is equivalent to $t\ge\max(u-x,x-n+1)$, so the inner sum is
$w_{\mathrm I}(x)=\#\{t\in[p_0,p_1]: t\ge\max(u-x,x-n+1)\}$, which is the displayed
expression. The function $x\mapsto\max(u-x,x-n+1,p_0)$ is convex, hence
$x\mapsto p_1-\max(\cdots)+1$ is concave, and $w_{\mathrm I}$ is its positive part, so
$w_{\mathrm I}\in\PFtwo$ by Lemma~\ref{lem:trunc}. An interval indicator is $\PFtwo$, so
$G_{\mathrm I}\in\PFtwo$ by (P3).

In region III the second interval $[c_t,d_t]=[C,D]$ is constant and $t\le x\le u-1-t$ is
equivalent to $t\le\min(x,u-1-x)$; the same argument applies, now with
$x\mapsto\min(x,u-1-x,q_1)$ concave.
\end{proof}

\begin{proposition}[regions II and IV: a constant support]\label{prop:regII}
In region II one has $a_t+c_t=u$ and $b_t+d_t=B+D$ for every $t$; in region IV one has
$a_t+c_t=A+C$ and $b_t+d_t=u+n-2$ for every $t$. Consequently, writing
$(\ell,h)=(u,B+D)$ for region II and $(\ell,h)=(A+C,u+n-2)$ for region IV, and
$H(s)=\min(s-\ell+1,\ h-s+1)$,
\[
  G_{\mathrm R}(s)=\rho_{\mathrm R}\bigl(H(s)\bigr),\qquad
  \rho_{\mathrm R}(v)=\sum_{t\in T_{\mathrm R}}\min\bigl(v,\ n_t,\ m_t\bigr)_+ ,
\]
where $n_t=b_t-a_t+1$, $m_t=d_t-c_t+1$. Both $G_{\mathrm{II}}$ and $G_{\mathrm{IV}}$ are
\emph{concave} on their support, hence $\PFtwo$.
\end{proposition}
\begin{proof}
Region II: $a_t+c_t=t+(u-t)=u$ and $b_t+d_t=B+D$. Region IV:
$a_t+c_t=A+C$ and $b_t+d_t=(u-1-t)+(t+n-1)=u+n-2$. In either case all the trapezoids
$\mathbf 1_{[a_t,b_t]}*\mathbf 1_{[c_t,d_t]}$ have the \emph{same} support $[\ell,h]$, and
the convolution of two interval indicators of widths $n_t,m_t$ is
\[
 \bigl(\mathbf 1_{[a_t,b_t]}*\mathbf 1_{[c_t,d_t]}\bigr)(s)
 =\min\bigl(s-\ell+1,\ h-s+1,\ n_t,\ m_t\bigr)_+
 =\min\bigl(H(s),\ n_t,\ m_t\bigr)_+ .
\]
Summing over $t\in T_{\mathrm R}$ gives $G_{\mathrm R}(s)=\rho_{\mathrm R}(H(s))$, using
$\min(\alpha,\beta)_+=\min(\alpha_+,\beta_+)$.

Now $v\mapsto\min(v,n_t,m_t)_+$ is nondecreasing and concave on $v\ge0$, so
$\rho_{\mathrm R}$ is nondecreasing and concave on $[0,\infty)$ with $\rho_{\mathrm R}(0)=0$.
The function $H$ is concave, and $H(s)\ge1$ for $s\in[\ell,h]$. A nondecreasing concave
function composed with a concave function is concave, so $G_{\mathrm R}=\rho_{\mathrm R}\circ H$
is concave on $[\ell,h]$; it is nonnegative there and its support is an interval. Apply
Lemma~\ref{lem:conc}.
\end{proof}

\begin{theorem}\label{thm:regions}
At $m=2$, for every cylindric skew shape, every level $n$ and every $b$, the sequence
$G$ of~\eqref{eq:G} decomposes as
$G=G_{\mathrm I}+G_{\mathrm{II}}+G_{\mathrm{III}}+G_{\mathrm{IV}}$ with each nonzero
summand in $\PFtwo$, by Propositions~\ref{prop:regI} and~\ref{prop:regII}.
\end{theorem}

\begin{corollary}[single-region criterion]\label{cor:single}
Let $T=[T_-,T_+]$ be the set of $t$ for which both intervals in~\eqref{eq:G} are nonempty.
If
\[
 \bigl(T_+<\tau_1\ \text{and}\ T_+\le\tau_2\bigr)\ \text{ or }\
 \bigl(T_-\ge\tau_1\ \text{and}\ T_+\le\tau_2\bigr)\ \text{ or }\
 \bigl(T_-\ge\tau_1\ \text{and}\ T_->\tau_2\bigr)\ \text{ or }\
 \bigl(T_+<\tau_1\ \text{and}\ T_->\tau_2\bigr)
\]
--- that is, if $T$ lies in a single one of the four cells --- then $k(\cdot,b)$ is $\PFtwo$;
in particular condition~(A) holds for this $(\lambda/\mu,n,b)$.
\end{corollary}

In a sweep over $2\le n\le9$, $m=2$, $d\le12$ ($5187$ nonempty slices) the criterion of
Corollary~\ref{cor:single} is satisfied by $4147$ slices, i.e.\ $79.9\%$.

\subsection{The gap}\label{sec:gap}

What remains is exactly the mixed terms. By the symmetrisation identity recalled
in~\eqref{eq:star}, Theorem~\ref{thm:regions} yields log-concavity of $G$ as soon as
\eqref{eq:star} holds for every pair of \emph{distinct} regions
$\mathrm R\ne\mathrm R'\in\{\mathrm I,\mathrm{II},\mathrm{III},\mathrm{IV}\}$:
\begin{equation}\label{eq:gap}
 2\,G_{\mathrm R}(s)G_{\mathrm R'}(s)\ \ge\
 G_{\mathrm R}(s-1)G_{\mathrm R'}(s+1)+G_{\mathrm R}(s+1)G_{\mathrm R'}(s-1)
 \qquad\text{for all }s.
\end{equation}

\textbf{Status of~\eqref{eq:gap}: verified, not proved.} Over $2\le n\le8$, $m=2$,
$d\le10$ ($2655$ slices), \eqref{eq:gap} holds on all $525$ occurring region pairs:
$\mathrm{I}\!-\!\mathrm{II}$ $177/177$, $\mathrm{II}\!-\!\mathrm{III}$ $177/177$,
$\mathrm{I}\!-\!\mathrm{III}$ $81/81$, $\mathrm{I}\!-\!\mathrm{IV}$ $45/45$,
$\mathrm{III}\!-\!\mathrm{IV}$ $45/45$.

Two structural facts constrain any proof of~\eqref{eq:gap}. First, all four supports are
nested inside the common maximal support: at $t=\tau_1$ the region I support
$[A+u-t,\ B+n-1+t]$ equals $[u,B+D]$, which is exactly the region II support, and at
$t=\tau_2$ the region III support $[t+C,\ u-1-t+D]$ also equals $[u,B+D]$. Second,
\eqref{eq:gap} is a statement about a \emph{cancellation across the regions}: it is
\emph{not} implied by any slice-local comparison, since \eqref{eq:star} already fails
between individual slice members (Section 2.1, item 1).

\section{Concavity is false}\label{sec:conc}

Proposition~\ref{prop:regII} proves more than log-concavity for regions II and IV: it
proves \emph{concavity}. This cannot be the general mechanism.

\begin{proposition}\label{prop:notconc}
The sequence $a\mapsto k(a,b)$ need not be concave on its support, already at $m=2$.
\end{proposition}
\begin{proof}
At $n=8$, $m=2$, $\mu=(0,3)$, $\lambda=(4,7)$, $b=2$ one finds
$k(\cdot,2)=(1,3,6,7,6,3,1)$; concavity fails at the second entry, where
$2\cdot3=6<1+6=7$. (Log-concavity holds there: $3^2=9\ge1\cdot6=6$.)
At $n=7$, $m=3$, $\mu=(0,1,4)$, $\lambda=(2,3,6)$, $b=2$ one finds $k(\cdot,2)=(1,3,6,3,1)$,
and again $2\cdot3=6<1+6=7$.
\end{proof}

Over a sweep with $2\le n\le9$, $m\le7$, $d\le11$ ($96781$ interior points), concavity
fails at $1758$ points --- distributed over $m=2$ ($44$), $m=3$ ($964$), $m=4$ ($720$),
$m=5$ ($30$) --- while log-concavity fails at $0$. So concavity is a small-case artifact
and log-concavity is the correct statement; in particular the concave-composition argument
of Proposition~\ref{prop:regII} \emph{must} break when more than one region is active,
which is consistent with the gap being exactly~\eqref{eq:gap}.

\section{What $m\ge3$ needs}

Proposition~\ref{prop:reform} is valid for every $m$. The special feature of $m=2$ is that
the box-slice $\{\nu\in\Bx(\mu):S_\nu=u\}$ is one-dimensional, so the whole family is
indexed by a single parameter $t$ and the four endpoint functions have only two
breakpoints. For $m\ge3$ the slice has dimension $m-1$ and no such parametrisation exists;
neither the reduction of Proposition~\ref{prop:m2red} nor the region decomposition of
Theorem~\ref{thm:regions} is available. We record this as the precise scope of the present
work rather than as a technical remark: the results of Section~3 are statements about
$m=2$ and nothing in their proofs suggests an $m$-independent form.

\section{Brändén--Huh: what the closure theorems do and do not give}

The programme had recorded ``read Brändén--Huh \texttt{1902.03719} at source'' as its
highest-value next step. This was done. The findings relevant here, cited by internal
\LaTeX\ label (an e-print theorem \emph{number} is a fact about a build we do not control):

\begin{itemize}
\item Thm~\verb|\label{flow}|: if $f\in\mathrm L^d_n$ then $f(Av)\in\mathrm L^d_m$ for any
nonnegative $n\times m$ matrix $A$ --- no positivity, invertibility or squareness.
\item Cor~\verb|\label{derivatives}|: $\sum_i a_i\partial_i f\in\mathrm L^{d-1}_n$ for
$a_i\ge0$.
\item Cor~\verb|\label{CorollaryProduct}|: \textbf{Lorentzian polynomials are closed under
products}, and Cor~\verb|\label{CorollaryConvolution}|: $N(f),N(g)$ Lorentzian
$\Rightarrow N(fg)$ Lorentzian. This was \emph{not} known to the programme and the brief
explicitly disclaimed it.
\item Closure under \emph{addition} is never addressed in the paper; it is false, and the
registry already holds the witness ($({x+y})^2+(2x+y)^2$).
\item There is \emph{no} characterisation of Lorentzian polynomials as limits of volume
polynomials; that was Question~\verb|\label{RealizationQuestion}|, answered negatively in
an added-in-proof footnote. So the ``realise $N(s^c)$ as a volume polynomial'' route has no
theorem behind it.
\item Cor~\verb|\label{normalizedcoefficients}|: if $\nu_f(\alpha)=\log c_\alpha$ is an
M-concave function then $f=\sum(c_\alpha/\alpha!)w^\alpha$ is Lorentzian. Sufficient, not
necessary.
\item The paper proves nothing whatever about Schur polynomials.
\end{itemize}

Two consequences for this programme. \emph{Negative:} none of these shortcuts
condition~(A). One might hope to obtain $h_b(x,z)=[y^b]s^c(x,y,z)$ by applying
$\partial_y^b$ (Cor~\verb|derivatives|) and then $y=0$ (the $\theta=0$ case of
Thm~\verb|flow|) to $N(s^c)$ --- but that presupposes $N(s^c)$ Lorentzian, which is the
whole conjecture. Condition (A) is a strictly weaker necessary condition and must be proved
on its own. \emph{Positive:} Cor~\verb|normalizedcoefficients| gives a genuinely new route
to the \emph{main} conjecture, since (L2) --- M-convexity of $\supp K^c$ --- is already
proved, so the local form of M-concavity (only $|\alpha-\beta|_1=4$ needs checking) applies.
At $\ell=3$ the case $\alpha-\beta=(2,-2,0)$ of that local axiom \emph{is} exactly
condition~(A). This is recorded as a new node, not pursued here.

\section{Verification}

All arithmetic is exact integer arithmetic. Two independent implementations of $K^c$ at
$\ell=3$ were written for this session --- the box model of Proposition~\ref{prop:reform}
and a direct enumeration of chains $\mu\prec\kappa\prec\sigma\prec\lambda$ --- and agree on
$647/647$ instances ($2\le n\le6$, $1\le m\le n$, $d\le6$).

\begin{center}
\begin{tabular}{lll}
\hline
claim & scope & result\\
\hline
$k$ via box model $=$ $k$ via chains & $n\le6,m\le n,d\le6$ & $647/647$\\
Proposition~\ref{prop:reform} reproduces $k$ & $n\le5,m\le n,d\le6$ & $227/227$\\
condition (A) & $n\le9,m\le7,d\le11$ & $96781/96781$ interior points\\
Proposition~\ref{prop:m2red} exact & $n\le9,m=2,d\le11$ & $6464$ slices, exact\\
$\Phi$ concave, $\Psi$ convex, both 1-Lipschitz & same & $0$ violations\\
Conjecture~\ref{conj:T}, exhaustive & profiles len $\le6$, shifts $[-4,6]$ & $29000/29000$\\
Conjecture~\ref{conj:T}, random & slopes $\{-1,0,1\}$ & $29251/29251$\\
four-region closed forms & $n\le9,m=2,d\le12$ & I $1363/1363$, II $2988/2988$,\\
 & & III $1363/1363$, IV $579/579$\\
$G=\sum_{\mathrm R}G_{\mathrm R}$ & $n\le8,m=2,d\le10$ & $2655/2655$\\
each nonempty region $\PFtwo$ & same & $0$ failures\\
single-region criterion correct & $n\le9,m=2,d\le12$ & $5187/5187$; $4147$ single\\
inter-region \eqref{eq:gap} & same & $525/525$ \textbf{(gap)}\\
\hline
\multicolumn{3}{l}{\textsc{negative controls} --- each bites, so the hypotheses are load-bearing}\\
\hline
Conjecture~\ref{conj:T} with slopes $\pm2$ & random & $117$ failures / $32900$\\
Conjecture~\ref{conj:T}, no slope bound & random & $40$ failures / $19133$\\
pairwise \eqref{eq:star} inside a slice & $n\le6,m\ge2,d\le7$ & $4$ failures / $9871$\\
general difference-constraint polytope & random, $m=2$ & $312$ failures / $6474$\\
concavity of $k(\cdot,b)$ & $n\le9,m\le7,d\le11$ & $1758$ failures / $96781$\\
$f_b$ real-rooted & $n\le6,m\le n,d\le7$ & $600$ failures / $1258$\\
$f_b$ product of cyclotomics & same & $91$ failures / $1176$\\
greedy induction on $t$ & random & fails $47\%$\\
greedy induction on columns & random & fails $69\%$\\
\hline
\end{tabular}
\end{center}

\section{Summary}

\begin{itemize}
\item \textbf{Proved, all $m$:} condition (A) is equivalent to log-concavity of a sum of
$\PFtwo$ sequences over one slice of a box (Proposition~\ref{prop:reform}).
\item \textbf{Proved, $m=2$:} an exact reduction to two $1$-Lipschitz profiles
(Proposition~\ref{prop:m2red}); a four-region decomposition with each region $\PFtwo$
(Theorem~\ref{thm:regions}); hence condition (A) whenever a single region is active
(Corollary~\ref{cor:single}), $79.9\%$ of slices in our sweep.
\item \textbf{Gap, precisely:} inequality~\eqref{eq:gap} between distinct regions.
Verified $525/525$; unproved.
\item \textbf{Open:} $m\ge3$ entirely, beyond Proposition~\ref{prop:reform}.
\item \textbf{Not claimed:} (A) does not settle (L3). By \texttt{l3-det-reduction},
(L3) at $\ell=3$ needs $e_2(M)\le0$ \emph{and} $\det M\ge0$; (A) gives only the former,
and the latter is condition (B).
\end{itemize}

\end{document}
